\section{Introduction}

The deployment of sequential decision-making algorithms, particularly Reinforcement Learning (RL), in mobile health (mHealth) and Ecological Momentary Assessment (EMA) platforms holds transformative potential for personalizing psychiatric interventions. By continuously adapting to incoming patient data, these algorithms can theoretically deliver the right intervention at the optimal moment. However, the operationalization of highly adaptive RL in live clinical trials introduces a fundamental tension between proximal algorithmic optimization and the preservation of rigorous scientific utility—a conflict we term the ``Adaptivity versus Fidelity'' trade-off.

As established by Gottesman et al. (2019), ``evaluation is harder than optimization'' in the context of clinical RL; achieving high cumulative rewards in offline training provides no guarantee of prospective safety or clinical generalizability. In live, dynamic trials, unconstrained RL agents often exhibit a phenomenon where they ``game'' the clinical environment. Driven exclusively by the mandate to maximize a predefined proximal reward signal, these agents frequently converge upon hyper-specific intervention types or narrow temporal niches. While this over-exploitation may superficially inflate early engagement metrics, it effectively eliminates the policy stochasticity necessary to compute reliable counterfactual estimators. Consequently, this destruction of data variance violates the fundamental requirements for Off-Policy Evaluation (OPE), rendering the trial incapable of producing generalized causal inferences or statistically valid comparisons between treatment arms at its conclusion.

Furthermore, ensuring algorithm fidelity—the degree to which an adaptive system adheres to the safety and exploratory protocols of its original clinical design—is paramount in psychiatric contexts where patient states are highly volatile (Trella et al., 2024). Traditional static safety constraints offer a baseline level of safeguarding but are inherently inflexible; they fail to capitalize on periods of clinical stability (euthymia) where broader exploration could safely generate the high-variance causal data essential for robust OPE. Conversely, unchecked algorithmic adaptivity risks exposing patients to unvalidated, edge-case intervention sequences during periods of acute psychological vulnerability.

To resolve this clinical and statistical impasse, we introduce a Dynamic Fidelity-Constrained RL architecture. This framework explicitly counterbalances reward optimization with the real-time continuous estimation of statistical power and algorithmic fidelity. By dynamically modulating the agent's exploration boundaries based on the immediate latent clinical state and the accumulating capacity for causal evaluation, our approach ensures both the safety of the participant and the indispensable scientific integrity of the trial. This is operationalized via Kullback-Leibler (KL) divergence penalties, drawing on the principles of Behavior Regularized Actor Critic (BRAC) (Wu et al., 2019) to anchor exploratory policies to safe, historically validated clinical baselines. Furthermore, our approach functions as a form of Distributionally Robust Optimization (DRO) (arXiv:2024), providing theoretical safety buffers against non-stationary distribution shifts.
